\documentclass[12pt,a4paper,openany,oneside]{book}

% --- Paquetes ---
\usepackage[utf8]{inputenc}
\usepackage[spanish, es-tabla]{babel}
\usepackage{amsmath, amsfonts, amssymb}
\usepackage{graphicx}
\usepackage{geometry}
\usepackage{hyperref}
\usepackage{booktabs}
\usepackage{fancyhdr}
\usepackage{listings}
\usepackage{xcolor}
\usepackage{tikz}
\usepackage{enumitem}
\usepackage{array}
\usetikzlibrary{arrows.meta, positioning, shapes.geometric}

\geometry{margin=2.5cm}

% --- Estilo de Código Python ---
\definecolor{codegreen}{rgb}{0,0.6,0}
\definecolor{codegray}{rgb}{0.5,0.5,0.5}
\definecolor{codepurple}{rgb}{0.58,0,0.82}
\definecolor{backcolour}{rgb}{0.95,0.95,0.92}

\lstset{
    language=Python,
    backgroundcolor=\color{backcolour},   
    commentstyle=\color{codegreen},
    keywordstyle=\color{blue},
    numberstyle=\tiny\color{codegray},
    stringstyle=\color{codepurple},
    basicstyle=\ttfamily\small,
    breaklines=true,
    frame=single,
    showstringspaces=false
}

% --- Cabecera y Pie de Página ---
\pagestyle{fancy}
\fancyhf{}
\renewcommand{\headrulewidth}{0.4pt}
\renewcommand{\footrulewidth}{0.4pt}
\fancyhead[L]{\small Carlos César Sánchez Coronel}
\fancyhead[R]{\small Machine Learning para IA}
\fancyfoot[L]{\small Carlos César Sánchez Coronel}
\fancyfoot[C]{\thepage}

% --- Portada ---
\title{
    \textbf{Sílabo del Curso} \\ 
    \Huge \textbf{MACHINE LEARNING} \\
    \Large \textsc{para Inteligencia Artificial} \\[1cm]
    \normalsize Fundamentos, algoritmos y aplicaciones prácticas
}
\author{
    \small Por \\ \vspace{0.2cm}
    \Large \textsc{Carlos César Sánchez Coronel}
}
\date{2026}

\begin{document}

\maketitle
\tableofcontents

% ------------------------------------------------------------------
% PRESENTACIÓN DEL CURSO
% ------------------------------------------------------------------
\chapter*{Presentación del Curso}
\addcontentsline{toc}{chapter}{Presentación del Curso}

Este curso está diseñado para futuros ingenieros e científicos de datos que deseen dominar los fundamentos del machine learning desde una perspectiva práctica y orientada a la industria. A lo largo de 14 semanas, exploraremos los algoritmos más importantes, técnicas de preprocesamiento, evaluación, interpretabilidad y despliegue, con un enfoque en las preguntas y habilidades que se esperan en entrevistas técnicas y en el trabajo diario.

Cada semana incluye:
\begin{itemize}
    \item \textbf{Logro de la sesión}: objetivo concreto.
    \item \textbf{Conceptos}: fundamentos teóricos.
    \item \textbf{Aplicaciones y casos típicos}: ejemplos reales de uso.
    \item \textbf{Métricas}: indicadores de rendimiento relevantes.
    \item \textbf{Python}: librerías y código de ejemplo.
\end{itemize}

\vspace{0.5cm}
\textbf{Estructura del curso (14 semanas):}
\begin{itemize}
    \item[$\square$] Semanas 1-3: Fundamentos y preprocesamiento
    \item[$\square$] Semanas 4-8: Modelos supervisados
    \item[$\square$] Semanas 9-11: Técnicas especializadas
    \item[$\square$] Semanas 12-14: Interpretabilidad y despliegue
\end{itemize}

% ------------------------------------------------------------------
% SEMANA 1: INTRODUCCIÓN Y CICLO DE VIDA
% ------------------------------------------------------------------
\chapter{Semana 1: Introducción a Machine Learning y Ciclo de Vida}

\section{Logro de la sesión}
Comprender el flujo completo de un proyecto de machine learning, identificar los tipos de aprendizaje y reconocer casos de uso en la industria.

\subsection{Conceptos}
\begin{itemize}
    \item Definición de Machine Learning y su relación con la IA.
    \item Tipos de aprendizaje: supervisado, no supervisado, por refuerzo y semi-supervisado.
    \item Ciclo de vida de un proyecto ML: definición del problema, recolección de datos, preprocesamiento, modelado, evaluación, despliegue y monitoreo.
    \item Roles involucrados: Data Scientist, ML Engineer, Data Analyst.
\end{itemize}

\subsection{Aplicaciones y casos típicos}
\begin{itemize}
    \item Sistemas de recomendación (Netflix, Amazon).
    \item Detección de fraude en transacciones bancarias.
    \item Mantenimiento predictivo en maquinaria industrial.
    \item Clasificación de imágenes (diagnóstico médico, vehículos autónomos).
    \item Predicción de series temporales (ventas, demanda energética).
\end{itemize}

\subsection{Métricas}
\begin{itemize}
    \item No aplican métricas específicas; se introducen conceptos de evaluación general.
\end{itemize}

\subsection{Python}
\begin{itemize}
    \item Introducción a scikit-learn: estructura de la librería, datasets de ejemplo (iris, boston).
    \item Carga y exploración básica con pandas.
\end{itemize}

% ------------------------------------------------------------------
% SEMANA 2: ANÁLISIS EXPLORATORIO Y FEATURE ENGINEERING
% ------------------------------------------------------------------
\chapter{Semana 2: Análisis Exploratorio y Feature Engineering}

\section{Logro de la sesión}
Realizar un análisis exploratorio eficaz y aplicar técnicas de feature engineering para preparar datos de calidad.

\subsection{Conceptos}
\begin{itemize}
    \item Análisis Exploratorio de Datos (EDA): estadísticas descriptivas, detección de outliers, análisis de correlaciones.
    \item Visualizaciones: histogramas, boxplots, pairplots, mapas de calor.
    \item Feature engineering: creación de nuevas características, transformaciones (log, polinómicas, interacciones), discretización.
    \item Manejo de datos faltantes y outliers.
    \item Pipelines de preprocesamiento.
\end{itemize}

\subsection{Aplicaciones y casos típicos}
\begin{itemize}
    \item Identificación de patrones en datos de clientes para segmentación.
    \item Creación de variables de tiempo (día de la semana, mes) en series temporales.
    \item Normalización de imágenes para modelos de visión.
\end{itemize}

\subsection{Métricas}
\begin{itemize}
    \item No se evalúan modelos, pero se usan métricas de correlación (Pearson, Spearman) y estadísticas descriptivas.
\end{itemize}

\subsection{Python}
\begin{itemize}
    \item pandas para manipulación y estadísticas.
    \item matplotlib y seaborn para visualizaciones.
    \item \texttt{sklearn.preprocessing} para escalado, codificación, imputación.
    \item \texttt{sklearn.pipeline.Pipeline}.
\end{itemize}

% ------------------------------------------------------------------
% SEMANA 3: REGRESIÓN LINEAL Y REGULARIZACIÓN
% ------------------------------------------------------------------
\chapter{Semana 3: Regresión Lineal y Regularización}

\section{Logro de la sesión}
Construir e interpretar modelos de regresión lineal, aplicando regularización para mejorar su generalización.

\subsection{Conceptos}
\begin{itemize}
    \item Regresión lineal simple y múltiple: formulación, interpretación de coeficientes.
    \item Supuestos del modelo: linealidad, independencia, homocedasticidad, normalidad de residuos.
    \item Regularización: Ridge (L2), Lasso (L1) y Elastic Net.
    \item Interpretación de la regularización y selección de hiperparámetros.
\end{itemize}

\subsection{Aplicaciones y casos típicos}
\begin{itemize}
    \item Predicción de precios de viviendas (Kaggle House Prices).
    \item Estimación de demanda de productos.
    \item Análisis de impacto de variables en marketing.
\end{itemize}

\subsection{Métricas}
\begin{itemize}
    \item Error Cuadrático Medio (MSE), Raíz del Error Cuadrático Medio (RMSE), Error Absoluto Medio (MAE), Coeficiente de determinación \(R^2\).
\end{itemize}

\subsection{Python}
\begin{itemize}
    \item \texttt{sklearn.linear\_model.LinearRegression}, \texttt{Ridge}, \texttt{Lasso}, \texttt{ElasticNet}.
    \item \texttt{sklearn.metrics} para cálculo de métricas.
    \item Validación cruzada con \texttt{GridSearchCV}.
\end{itemize}

% ------------------------------------------------------------------
% SEMANA 4: CLASIFICACIÓN I - REGRESIÓN LOGÍSTICA Y MÉTRICAS
% ------------------------------------------------------------------
\chapter{Semana 4: Clasificación I - Regresión Logística y Métricas}

\section{Logro de la sesión}
Implementar modelos de clasificación binaria y evaluarlos correctamente con las métricas adecuadas.

\subsection{Conceptos}
\begin{itemize}
    \item Regresión logística: función sigmoide, interpretación probabilística, límite de decisión.
    \item Matriz de confusión: verdaderos positivos, falsos positivos, etc.
    \item Métricas derivadas: precisión, recall (sensibilidad), especificidad, F1-score.
    \item Curva ROC y AUC (Area Under the Curve).
    \item Manejo de clases desbalanceadas: técnicas de muestreo (SMOTE, undersampling), pesos en la función de pérdida.
\end{itemize}

\subsection{Aplicaciones y casos típicos}
\begin{itemize}
    \item Detección de spam en correos electrónicos.
    \item Diagnóstico médico (enfermedad vs sano).
    \item Clasificación de clientes (compradores/no compradores).
\end{itemize}

\subsection{Métricas}
\begin{itemize}
    \item Precisión, recall, F1-score, AUC-ROC.
\end{itemize}

\subsection{Python}
\begin{itemize}
    \item \texttt{sklearn.linear\_model.LogisticRegression}.
    \item \texttt{sklearn.metrics}: \texttt{confusion\_matrix}, \texttt{classification\_report}, \texttt{roc\_curve}, \texttt{roc\_auc\_score}.
    \item \texttt{imbalanced-learn} para SMOTE.
\end{itemize}

% ------------------------------------------------------------------
% SEMANA 5: CLASIFICACIÓN II - ALGORITMOS CLÁSICOS
% ------------------------------------------------------------------
\chapter{Semana 5: Clasificación II - Algoritmos Clásicos}

\section{Logro de la sesión}
Conocer y aplicar algoritmos de clasificación clásicos como KNN, Naive Bayes y SVM, entendiendo sus fundamentos y limitaciones.

\subsection{Conceptos}
\begin{itemize}
    \item K-Vecinos más cercanos (KNN): distancia, elección de $k$, influencia de la escala.
    \item Naive Bayes: teorema de Bayes, supuesto de independencia, variantes (Gaussian, Multinomial, Bernoulli).
    \item Máquinas de Vectores de Soporte (SVM): margen máximo, kernel lineal y no lineal (RBF), parámetros $C$ y $\gamma$.
\end{itemize}

\subsection{Aplicaciones y casos típicos}
\begin{itemize}
    \item KNN: sistemas de recomendación (basados en similitud), reconocimiento de dígitos.
    \item Naive Bayes: clasificación de textos (filtro de spam, análisis de sentimiento).
    \item SVM: reconocimiento facial, clasificación de proteínas.
\end{itemize}

\subsection{Métricas}
\begin{itemize}
    \item Las mismas que en clasificación (precisión, recall, F1, AUC).
\end{itemize}

\subsection{Python}
\begin{itemize}
    \item \texttt{sklearn.neighbors.KNeighborsClassifier}.
    \item \texttt{sklearn.naive\_bayes.GaussianNB}, \texttt{MultinomialNB}, \texttt{BernoulliNB}.
    \item \texttt{sklearn.svm.SVC}.
\end{itemize}

% ------------------------------------------------------------------
% SEMANA 6: ÁRBOLES DE DECISIÓN Y RANDOM FOREST
% ------------------------------------------------------------------
\chapter{Semana 6: Árboles de Decisión y Random Forest}

\section{Logro de la sesión}
Construir e interpretar árboles de decisión y modelos de ensamble basados en bagging (Random Forest).

\subsection{Conceptos}
\begin{itemize}
    \item Árboles de clasificación y regresión: criterios de división (Gini, entropía, MSE), poda, sobreajuste.
    \item Interpretación de árboles: importancia de características, visualización.
    \item Bagging (Bootstrap Aggregating) y Random Forest: muestreo con reemplazo, aleatoriedad en la selección de características, Out-of-Bag (OOB) error.
\end{itemize}

\subsection{Aplicaciones y casos típicos}
\begin{itemize}
    \item Detección de fraudes en tarjetas de crédito.
    \item Análisis de riesgo crediticio.
    \item Problemas de alta dimensionalidad (genómica, texto).
\end{itemize}

\subsection{Métricas}
\begin{itemize}
    \item Para clasificación: precisión, recall, F1, AUC. Para regresión: MSE, RMSE, MAE, \(R^2\).
    \item Importancia de características (Gini importance, permutación).
\end{itemize}

\subsection{Python}
\begin{itemize}
    \item \texttt{sklearn.tree.DecisionTreeClassifier}, \texttt{DecisionTreeRegressor}.
    \item \texttt{sklearn.ensemble.RandomForestClassifier}, \texttt{RandomForestRegressor}.
    \item Visualización con \texttt{plot\_tree} y \texttt{export\_graphviz}.
\end{itemize}

% ------------------------------------------------------------------
% SEMANA 7: GRADIENT BOOSTING (XGBoost, LightGBM)
% ------------------------------------------------------------------
\chapter{Semana 7: Gradient Boosting (XGBoost, LightGBM)}

\section{Logro de la sesión}
Dominar los algoritmos de boosting modernos y ajustar sus hiperparámetros para obtener modelos de alto rendimiento.

\subsection{Conceptos}
\begin{itemize}
    \item Boosting: concepto de ensamblaje secuencial, corrección de errores.
    \item Gradient Boosting: función de pérdida, árboles débiles, learning rate.
    \item Implementaciones eficientes: XGBoost (regularización, manejo de nulos, early stopping), LightGBM (muestreo por gradiente, características categóricas).
    \item Comparación con Random Forest.
\end{itemize}

\subsection{Aplicaciones y casos típicos}
\begin{itemize}
    \item Competencias Kaggle (ganadoras de premios).
    \item Predicción de clics en publicidad online.
    \item Modelado de tablas en general.
\end{itemize}

\subsection{Métricas}
\begin{itemize}
    \item Dependen del problema: clasificación (log-loss, AUC) o regresión (RMSE, MAE).
\end{itemize}

\subsection{Python}
\begin{itemize}
    \item \texttt{xgboost.XGBClassifier}, \texttt{XGBRegressor}.
    \item \texttt{lightgbm.LGBMClassifier}, \texttt{LGBMRegressor}.
    \item Búsqueda de hiperparámetros con \texttt{GridSearchCV} o \texttt{RandomizedSearchCV}.
\end{itemize}

% ------------------------------------------------------------------
% SEMANA 8: EVALUACIÓN Y VALIDACIÓN DE MODELOS
% ------------------------------------------------------------------
\chapter{Semana 8: Evaluación y Validación de Modelos}

\section{Logro de la sesión}
Evaluar correctamente el rendimiento de los modelos, detectar overfitting y seleccionar los mejores hiperparámetros.

\subsection{Conceptos}
\begin{itemize}
    \item Underfitting y overfitting: causas y soluciones.
    \item Compromiso sesgo-varianza (Bias-Variance tradeoff).
    \item Validación cruzada: k-fold, stratified k-fold, leave-one-out.
    \item Curvas de aprendizaje: diagnóstico de under/overfitting.
    \item Búsqueda de hiperparámetros: Grid Search, Random Search, Bayesian Optimization.
\end{itemize}

\subsection{Aplicaciones y casos típicos}
\begin{itemize}
    \item Ajuste de modelos en cualquier proyecto para garantizar generalización.
    \item Selección de hiperparámetros en competencias.
\end{itemize}

\subsection{Métricas}
\begin{itemize}
    \item Según el tipo de modelo, se usan las métricas ya vistas (RMSE, AUC, etc.) en validación cruzada.
\end{itemize}

\subsection{Python}
\begin{itemize}
    \item \texttt{sklearn.model\_selection}: \texttt{cross\_val\_score}, \texttt{learning\_curve}, \texttt{GridSearchCV}, \texttt{RandomizedSearchCV}.
\end{itemize}

% ------------------------------------------------------------------
% SEMANA 9: APRENDIZAJE NO SUPERVISADO - CLUSTERING Y PCA
% ------------------------------------------------------------------
\chapter{Semana 9: Aprendizaje No Supervisado - Clustering y PCA}

\section{Logro de la sesión}
Aplicar técnicas de reducción de dimensionalidad y clustering para explorar datos no etiquetados.

\subsection{Conceptos}
\begin{itemize}
    \item PCA (Análisis de Componentes Principales): maximización de varianza, interpretación de componentes, scree plot.
    \item K-Means: algoritmo, inicialización, elección de $k$ (método del codo, silueta).
    \item Otros algoritmos de clustering: DBSCAN (concepto, parámetros eps y minPts), clustering jerárquico (dendrogramas).
    \item Métricas de evaluación: coeficiente de silueta, índice de Davies-Bouldin.
\end{itemize}

\subsection{Aplicaciones y casos típicos}
\begin{itemize}
    \item Segmentación de clientes en marketing.
    \item Compresión de datos (PCA).
    \item Detección de anomalías con DBSCAN.
\end{itemize}

\subsection{Métricas}
\begin{itemize}
    \item Silueta, Davies-Bouldin, inercia (para K-Means).
\end{itemize}

\subsection{Python}
\begin{itemize}
    \item \texttt{sklearn.decomposition.PCA}.
    \item \texttt{sklearn.cluster.KMeans}, \texttt{DBSCAN}.
    \item \texttt{sklearn.metrics.silhouette\_score}.
\end{itemize}

% ------------------------------------------------------------------
% SEMANA 10: SERIES DE TIEMPO
% ------------------------------------------------------------------
\chapter{Semana 10: Series de Tiempo}

\section{Logro de la sesión}
Modelar datos temporales utilizando técnicas estadísticas y de machine learning.

\subsection{Conceptos}
\begin{itemize}
    \item Componentes de una serie temporal: tendencia, estacionalidad, ciclos, ruido.
    \item Descomposición aditiva y multiplicativa.
    \item Feature engineering temporal: rezagos (lags), ventanas móviles, diferencias.
    \item Modelos clásicos: ARIMA, SARIMA.
    \item Modelos de machine learning para series: regresión con características temporales, XGBoost, Prophet.
\end{itemize}

\subsection{Aplicaciones y casos típicos}
\begin{itemize}
    \item Pronóstico de ventas en retail.
    \item Predicción de demanda energética.
    \item Análisis de tráfico web.
\end{itemize}

\subsection{Métricas}
\begin{itemize}
    \item MAE, MAPE, RMSE, MASE.
\end{itemize}

\subsection{Python}
\begin{itemize}
    \item pandas para manejo de fechas y creación de lags.
    \item \texttt{statsmodels} para ARIMA y descomposición.
    \item Prophet (librería de Facebook).
    \item Evaluación con sklearn métricas.
\end{itemize}

% ------------------------------------------------------------------
% SEMANA 11: MODELOS COMPLEMENTARIOS (REGRESIÓN, CLASIFICACIÓN, CLUSTERING)
% ------------------------------------------------------------------
\chapter{Semana 11: Modelos Complementarios}

\section{Logro de la sesión}
Conocer otros algoritmos y técnicas que pueden aparecer en entrevistas o en problemas específicos, y entender cuándo aplicarlos.

\subsection{Conceptos}
\begin{itemize}
    \item Regresión polinómica: extensión de la lineal, riesgo de overfitting.
    \item Regresión robusta: Huber, RANSAC (para datos con outliers).
    \item SVM en regresión: SVR.
    \item DBSCAN a fondo: densidad, ruido, parámetros.
    \item Clustering jerárquico: dendrogramas, aglomerativo vs divisivo.
    \item Métricas de clustering adicionales: índice de Rand ajustado, información mutua.
\end{itemize}

\subsection{Aplicaciones y casos típicos}
\begin{itemize}
    \item Regresión robusta: cuando hay outliers significativos (precios de vivienda atípicos).
    \item DBSCAN: detección de clusters de forma arbitraria, identificación de ruido.
    \item Clustering jerárquico: análisis filogenético, segmentación con interpretación jerárquica.
\end{itemize}

\subsection{Métricas}
\begin{itemize}
    \item Para regresión robusta: mismas métricas que regresión.
    \item Para clustering: silueta, Davies-Bouldin, Rand index.
\end{itemize}

\subsection{Python}
\begin{itemize}
    \item \texttt{sklearn.linear\_model.HuberRegressor}, \texttt{RANSACRegressor}.
    \item \texttt{sklearn.svm.SVR}.
    \item \texttt{sklearn.cluster.DBSCAN}, \texttt{AgglomerativeClustering}.
    \item \texttt{sklearn.metrics.adjusted\_rand\_score}.
\end{itemize}

% ------------------------------------------------------------------
% SEMANA 12: SISTEMAS DE RECOMENDACIÓN
% ------------------------------------------------------------------
\chapter{Semana 12: Sistemas de Recomendación}

\section{Logro de la sesión}
Diseñar e implementar sistemas de recomendación basados en filtrado colaborativo y factorización matricial.

\subsection{Conceptos}
\begin{itemize}
    \item Filtrado colaborativo basado en usuarios y en ítems.
    \item Factorización matricial: SVD (descomposición en valores singulares), Funk SVD.
    \item Métricas de evaluación: precisión@k, recall@k, NDCG, RMSE en predicciones de rating.
    \item Problemas: cold start, escalabilidad, sparse data.
\end{itemize}

\subsection{Aplicaciones y casos típicos}
\begin{itemize}
    \item Recomendación de productos en Amazon.
    \item Películas en Netflix.
    \item Música en Spotify.
\end{itemize}

\subsection{Métricas}
\begin{itemize}
    \item Para clasificación: precisión@k, recall@k, NDCG. Para predicción de rating: RMSE, MAE.
\end{itemize}

\subsection{Python}
\begin{itemize}
    \item Implementación con \texttt{surprise} (librería especializada) o usando SVD con \texttt{sklearn.decomposition.TruncatedSVD}.
    \item Alternativa: uso de embeddings con redes neuronales (opcional).
\end{itemize}

% ------------------------------------------------------------------
% SEMANA 13: INTERPRETABILIDAD DE MODELOS
% ------------------------------------------------------------------
\chapter{Semana 13: Interpretabilidad de Modelos}

\section{Logro de la sesión}
Explicar las predicciones de modelos complejos utilizando técnicas de interpretabilidad global y local.

\subsection{Conceptos}
\begin{itemize}
    \item Importancia de características en modelos basados en árboles.
    \item Partial Dependence Plots (PDP) e Individual Conditional Expectation (ICE).
    \item SHAP (SHapley Additive exPlanations): valores de Shapley, interpretación local y global.
    \item LIME (Local Interpretable Model-agnostic Explanations).
    \item Diferencias entre interpretabilidad global y local.
\end{itemize}

\subsection{Aplicaciones y casos típicos}
\begin{itemize}
    \item Explicar predicciones a stakeholders (por qué se negó un crédito).
    \item Cumplir regulaciones (GDPR, derecho a explicación).
    \item Depurar modelos (identificar sesgos).
\end{itemize}

\subsection{Métricas}
\begin{itemize}
    \item No hay métricas directas; se evalúa la consistencia de las explicaciones.
\end{itemize}

\subsection{Python}
\begin{itemize}
    \item \texttt{sklearn.inspection} para PDP y ICE.
    \item Librerías \texttt{shap} y \texttt{lime}.
\end{itemize}

% ------------------------------------------------------------------
% SEMANA 14: DESPLIEGUE DE MODELOS (MLOps)
% ------------------------------------------------------------------
\chapter{Semana 14: Despliegue de Modelos (MLOps)}

\section{Logro de la sesión}
Llevar un modelo a producción, entender los conceptos básicos de MLOps y monitorear su rendimiento.

\subsection{Conceptos}
\begin{itemize}
    \item Separación entre entrenamiento e inferencia.
    \item Serialización de modelos: pickle, joblib, ONNX.
    \item Creación de APIs con Flask o FastAPI.
    \item Contenerización básica con Docker.
    \item Versionamiento de modelos y datos (DVC, MLflow).
    \item Monitoreo de modelos en producción: data drift, concept drift, métricas de rendimiento.
\end{itemize}

\subsection{Aplicaciones y casos típicos}
\begin{itemize}
    \item Integrar modelo de recomendación en una app web.
    \item Actualizar modelos periódicamente en producción.
    \item Detectar degradación del modelo por cambios en los datos.
\end{itemize}

\subsection{Métricas}
\begin{itemize}
    \item Se monitorean las mismas métricas que en entrenamiento (RMSE, AUC, etc.) en datos nuevos.
    \item Detección de drift usando tests estadísticos (KS, PSI).
\end{itemize}

\subsection{Python}
\begin{itemize}
    \item Serialización con \texttt{pickle} y \texttt{joblib}.
    \item Ejemplo de API con Flask: carga del modelo y endpoint de predicción.
    \item Introducción a MLflow para registro de experimentos y versionado.
    \item Uso de \texttt{evidently} o \texttt{alibi-detect} para monitoreo.
\end{itemize}

% ------------------------------------------------------------------
% METODOLOGÍA Y EVALUACIÓN
% ------------------------------------------------------------------
\chapter{Metodología y Evaluación}

\section{Estrategias metodológicas}
\begin{itemize}
    \item \textbf{Clases síncronas:} Exposición conceptual con diapositivas, resolución de problemas en pizarra y programación en vivo con Python utilizando cuadernos de Jupyter.
    \item \textbf{Trabajo asíncrono:} Lectura de material complementario, resolución de ejercicios prácticos y pequeños proyectos semanales.
    \item \textbf{Aprendizaje activo:} Discusión de casos reales, ejercicios de entrevistas técnicas y retos de código.
    \item \textbf{Recursos:} Plataforma virtual, Zoom, grabaciones, foros y cuadernos de Jupyter.
\end{itemize}

\section{Materiales educativos}
\begin{itemize}
    \item Presentaciones interactivas (PPT, videos).
    \item Cuadernos de Jupyter con ejemplos y ejercicios.
    \item Bibliografía: libros clásicos (Géron, Hastie) y documentación oficial.
    \item Google Colab para ejecución sin instalación.
\end{itemize}

\section{Evaluación}
\begin{table}[h]
\centering
\begin{tabular}{lc}
\toprule
\textbf{Ítem} & \textbf{Ponderación} \\
\midrule
Evaluaciones continuas (EC) & 40\% \\
Examen parcial (EP) & 30\% \\
Examen final (EF) & 30\% \\
\midrule
\textbf{Promedio final (PF)} & PF = 0.4 \times EC + 0.3 \times EP + 0.3 \times EF \\
\bottomrule
\end{tabular}
\caption{Sistema de evaluación}
\end{table}
\begin{itemize}
    \item \textbf{Evaluaciones continuas:} Controles de lectura, entrega de ejercicios de programación, participación.
    \item \textbf{Examen parcial:} Cubre semanas 1-7.
    \item \textbf{Examen final:} Cubre semanas 8-14.
    \item \textbf{Examen sustitutorio:} Semana 15, reemplaza la nota más baja entre parcial y final.
\end{itemize}

% ------------------------------------------------------------------
% BIBLIOGRAFÍA
% ------------------------------------------------------------------
\chapter*{Bibliografía}
\addcontentsline{toc}{chapter}{Bibliografía}

\begin{itemize}
    \item Géron, A. (2019). \textit{Hands-On Machine Learning with Scikit-Learn, Keras, and TensorFlow}. O'Reilly.
    \item Hastie, T., Tibshirani, R., Friedman, J. (2009). \textit{The Elements of Statistical Learning}. Springer.
    \item Bishop, C. M. (2006). \textit{Pattern Recognition and Machine Learning}. Springer.
    \item VanderPlas, J. (2016). \textit{Python Data Science Handbook}. O'Reilly.
    \item Documentación oficial de scikit-learn, XGBoost, LightGBM, SHAP, etc.
\end{itemize}

\vspace{1cm}
\begin{flushright}
\textit{Elaborado por:} Mg. Carlos César Sánchez Coronel \\
\textit{Aprobado por:} Dirección de Ingeniería de Inteligencia Artificial \\
\textit{Fecha:} 29/03/2024
\end{flushright}

\end{document}